\documentclass[spanish,letterpaper,oneside]{article}

\def\documenttitle {Tipos de Aprendizaje en las Sociedades del Siglo XXI}
\def\documentsubtitle {Fundamentos para la enseñanza y el aprendizaje I · Cohorte SD26}
\def\documentsubject {Actividad 10 · Bloque I Tema 4 (Complemento Cynthia Luna Scott, 2015b)}
\def\documentauthor {Mtro. Martín Jonathan de la Cruz Muñoz}
\def\coursename {Fundamentos para la enseñanza y el aprendizaje I}
\def\coursecode {MAEM-SD26}

\def\universityname {Instituto de Investigación, Innovación y Estudios de Posgrado para la Educación del Estado de Nuevo León}
\def\universityfaculty {Maestría en Aprendizaje y Enseñanza de las Matemáticas}
\def\universitydepartment {Fundamentos para la enseñanza y el aprendizaje I}
\def\universitydepartmentimage {img/iiiepe/marca-iiiepe.png}
\def\universitydepartmentimagecfg {height=1.75cm}
\def\universitylocation {Monterrey, Nuevo León}
\def\studentfullname {Mtro. Martín Jonathan de la Cruz Muñoz}
\def\studentid {00874}
\def\studentemail {martin.delacruz@campus.iiiepe.edu.mx}
\def\programperiod {2do periodo, septiembre--diciembre de 2026 (SD26)}
\def\programmodality {Modalidad presencial}
\def\programshort {MAEM}

\def\authortable {
  \begin{tabular}{ll}
    \textbf{Alumno:} & \studentfullname \\
    Matrícula: & \studentid \\
    Correo institucional: & \studentemail \\
    Unidad: & \coursename \\
    Programa: & \programshort \\
    Periodo: & \programperiod \\
    & \\
    \multicolumn{2}{l}{Fecha de realización: 19 de noviembre de 2026} \\
    \multicolumn{2}{l}{\universitylocation}
  \end{tabular}
}

\input{template}

\usepackage{helvet}
\renewcommand{\familydefault}{\sfdefault}
\usepackage{array}
\usepackage{longtable}
\usepackage{pdflscape}
\usepackage{tikz}
\usepackage{pgfplots}
\pgfplotsset{compat=1.18}
\usepackage[most]{tcolorbox}
\usetikzlibrary{arrows.meta,positioning,calc,fit,backgrounds,shapes.geometric}

\definecolor{iiiepeNavy}{HTML}{163B5C}
\definecolor{iiiepeBlue}{HTML}{1F6F9F}
\definecolor{iiiepeAqua}{HTML}{20B7A8}
\definecolor{iiiepeOrange}{HTML}{F28C28}
\definecolor{iiiepePaper}{HTML}{F6F8F7}
\definecolor{iiiepeInk}{HTML}{1E2A32}
\definecolor{iiiepePale}{HTML}{EEF5F8}
\definecolor{iiiepeMint}{HTML}{E7F6F4}
\definecolor{iiiepeWarm}{HTML}{FFF3E6}
\def\portraittitlecolor {iiiepeNavy}
\setcitestyle{authoryear,open={(},close={)}}
\makeatletter
\renewcommand{\@biblabel}[1]{}
\makeatother
\setlength{\bibhang}{1.27cm}
\setlength{\bibsep}{6pt}
\AtBeginDocument{\renewcommand{\refname}{Referencias}}

\newcommand{\insertcoverwatermark}{%
  \AddToShipoutPictureBG*{%
    \begin{tikzpicture}[remember picture,overlay]
      \node[opacity=0.04,inner sep=0pt] at (current page.center) {%
        \includegraphics[width=0.65\paperwidth]{img/iiiepe/marca-iiiepe.png}%
      };
    \end{tikzpicture}%
  }%
}

\setlength{\parskip}{5pt}\setlength{\parindent}{0pt}
\begin{document}\insertcoverwatermark\templatePortrait\templatePagecfg\templateFinalcfg\fancyfoot{}\fancyfoot[L]{\scriptsize IIIEPE · MAEM · SD26}\fancyfoot[C]{\scriptsize\thepage}\fancyfoot[R]{\scriptsize Actividad 10 · Aprendizaje Siglo XXI}\setstretch{1.3}\raggedbottom
\section{Introducción y postura central}
El segundo informe de \citet{scott2015futuro2}, \textit{El futuro del aprendizaje 2: ¿Qué tipo de aprendizaje se necesita en el siglo XXI?}, traslada el análisis desde las razones estructurales del cambio educativo hacia la definición precisa de las competencias y modalidades de aprendizaje que demanda nuestra época. Mientras las rutinas mecánicas pueden ser automatizadas o estandarizadas, la ciudadanía democrática y el trabajo técnico calificado exigen sujetos capaces de aprender permanentemente, colaborar en redes plurales, comunicar argumentos fundamentados y resolver problemas inéditos con criterio ético \citep{delors1996educacion,unesco2015replantear}.

Mi postura es que clasificar competencias genéricas como «habilidades blandas» separadas de las disciplinas empobrece la enseñanza matemática; por el contrario, los tipos de aprendizaje descritos por \citet{scott2015futuro2} solo cobran vigencia formativa cuando se articulan con la modelación, la deducción lógica y el análisis cuantitativo en entornos reales. Además, desde mi práctica docente en CONALEP, CECyTE e ingeniería, sostengo que reorganizar los cuatro pilares educativos en secuencias didácticas claras permite vincular la precisión instrumental con el desarrollo humano integral. Es decir, no obstante la urgencia de la cobertura curricular técnica, comprender y transferir métodos exige tiempo reflexivo; por ello, cada pilar orienta decisiones concretas en el aula.

\section{Articulación de los cuatro pilares y competencias del siglo XXI}
Recuperando el marco clásico de la Comisión Internacional sobre la Educación para el Siglo XXI \citep{delors1996educacion}, \citet{scott2015futuro2} organiza las finalidades educativas contemporáneas en cuatro pilares interdependientes, enriquecidos con capacidades cognitivas de orden superior y competencias socioemocionales y metacognitivas. La Figura~\ref{fig:act10-pilares} esquematiza su convergencia dentro de la enseñanza matemática.

\begin{figure}[htbp]
\centering
\begin{tikzpicture}[font=\scriptsize, >=Latex]
  \node[draw=iiiepeNavy,fill=iiiepeNavy,text=white,rounded corners=5pt,text width=4.0cm,minimum height=1.55cm,align=center,inner sep=6pt] (center) at (0,0)
    {\textbf{\small APRENDIZAJES DEL SIGLO XXI}\\Competencia matemática, metacognición y ciudadanía};

  \node[draw=iiiepeBlue,fill=iiiepePale,thick,rounded corners=4pt,text width=5.1cm,align=left,inner sep=6pt] (p1) at (-4.3,2.55)
    {\textbf{\color{iiiepeBlue}1. Aprender a conocer}\\Comprensión relacional, indagación científica, lectura crítica de datos y hábitos de aprendizaje continuo \citep{scott2015futuro2}.};

  \node[draw=iiiepeAqua!80!black,fill=iiiepeMint,thick,rounded corners=4pt,text width=5.1cm,align=left,inner sep=6pt] (p2) at (4.3,2.55)
    {\textbf{\color{iiiepeAqua!75!black}2. Aprender a hacer}\\Modelación de sistemas, transferencia entre contextos, manejo ético de instrumentos y trabajo en equipo \citep{conalep2025mecatronica}.};

  \node[draw=iiiepeOrange!90!black,fill=iiiepeWarm,thick,rounded corners=4pt,text width=5.1cm,align=left,inner sep=6pt] (p3) at (-4.3,-2.55)
    {\textbf{\color{iiiepeOrange!85!black}3. Aprender a ser}\\Autonomía intelectual, perseverancia ante el bloqueo, reflexión metacognitiva y responsabilidad ética \citep{schunk2012teorias}.};

  \node[draw=iiiepeNavy,fill=iiiepePaper,thick,rounded corners=4pt,text width=5.1cm,align=left,inner sep=6pt] (p4) at (4.3,-2.55)
    {\textbf{\color{iiiepeNavy}4. Aprender a vivir juntos}\\Escucha respetuosa de argumentos divergentes, colaboración intercultural y bien común escolar \citep{unesco2015replantear}.};

  \draw[->,very thick,iiiepeBlue] (p1.south east) -- (center.north west);
  \draw[->,very thick,iiiepeAqua!80!black] (p2.south west) -- (center.north east);
  \draw[->,very thick,iiiepeOrange!85!black] (p3.north east) -- (center.south west);
  \draw[->,very thick,iiiepeNavy] (p4.north west) -- (center.south east);
\end{tikzpicture}
\caption{Integración de los cuatro pilares de Delors y Cynthia Luna Scott (2015b) en la formación matemática.}
\label{fig:act10-pilares}
\end{figure}

El cuadro siguiente desglosa los tipos de aprendizaje identificados en \citet{scott2015futuro2} y precisa su traducción metodológica para el trabajo cotidiano en matemáticas técnicas e ingeniería:

{\small\setstretch{1.15}\renewcommand{\arraystretch}{1.32}
\begin{longtable}{|>{\raggedright\arraybackslash}p{0.21\linewidth}|>{\raggedright\arraybackslash}p{0.35\linewidth}|>{\raggedright\arraybackslash}p{0.38\linewidth}|}\hline
\rowcolor{iiiepePale} \textbf{Pilar y competencias} & \textbf{Alcance en Scott (2015b)} & \textbf{Concreción en aula de Matemáticas} \\\hline
\endfirsthead
\hline
\rowcolor{iiiepePale} \textbf{Pilar y competencias} & \textbf{Alcance en Scott (2015b)} & \textbf{Concreción en aula de Matemáticas} \\\hline
\endhead
Aprender a conocer (pensamiento crítico y resolución compleja) & Dominio de conceptos fundamentales articulado con estrategias heurísticas para evaluar evidencias e información plural. & Distinguir cuándo una relación empírica sigue proporcionalidad directa o función afín verificando razones incrementales y unidades \citep{sawyer2008optimizing}. \\\hline
Aprender a hacer (creatividad, innovación y alfabetización digital) & Capacidad para diseñar prototipos, usar tecnología productiva y adaptar procedimientos ante condiciones variables. & Programar hojas de cálculo para simular tolerancias de medición contrastando siempre la salida gráfica contra el cálculo manual. \\\hline
Aprender a ser (metacognición, autonomía y autorregulación) & Aprender a aprender, regular la atención, gestionar la frustración ante problemas difíciles y ejercer juicio ético. & Incorporar bitácoras de corrección de errores donde cada alumno explique qué supuesto falló antes de corregir su ecuación \citep{eef2021feedback}. \\\hline
Aprender a vivir juntos (colaboración, comunicación y ciudadanía) & Cooperación respetuosa, negociación de perspectivas y compromiso con problemas reales de la comunidad. & Organizar mesas de revisión entre pares donde equipos defienden decisiones técnicas frente a restricciones reales de costo y seguridad \citep{sep2023pensamiento}. \\\hline
\end{longtable}}

\section{Aplicación reflexiva a la práctica docente en matemáticas}
Al reflexionar sobre cómo aplicar estos tipos de aprendizaje en mi práctica, reconozco que tradicionalmente los cursos técnicos concentran la atención en el pilar del \textit{aprender a hacer} reducido a ejecución procedimental. Sin embargo, formar técnicos e ingenieros para entornos cambiantes exige que el hacer esté iluminado por el conocer profundo, regulado por la autoconciencia metacognitiva del ser y validado éticamente en el diálogo del vivir juntos \citep{scott2015futuro2,pimienta2012estrategias}.

Para ejemplificar esta integración mediante un diseño \textbf{por realizar}, considérese una situación sobre control estadístico y variación de consumo energético en un taller escolar: dado un equipo cuya demanda basal de vacío es $P_0 = 120\text{ W}$ y su consumo adicional por pieza procesada es $45\text{ W/pieza}$, los estudiantes modelan $P(n) = 120 + 45n$, evalúan la viabilidad eléctrica del circuito alimentador ante $n = 25$ piezas ($P(25) = 1245\text{ W}$) y debaten en equipo alternativas para reducir pérdidas sin sacrificar calidad. De esta manera, el cálculo algebraico sirve al criterio técnico, a la reflexión metacognitiva sobre hipótesis límite y a la responsabilidad compartida \citep{sep2023pensamiento}. No obstante su pertinencia teórica, esta propuesta requiere verificación futura con evidencias reales de aula para calibrar ritmos y andamiajes diferenciados.

\clearpage
\section{Conclusiones}
En conclusión, el análisis de \citet{scott2015futuro2} ratifica que la educación del siglo XXI no demanda abandonar los conocimientos disciplinares sólidos, sino transformarlos en instrumentos vivos de pensamiento crítico, cooperación e innovación responsable. Considero que articular los cuatro pilares de \citet{delors1996educacion} con problemas matemáticos auténticos evita tanto el reduccionismo algorítmico como la retórica vacía de competencias descontextualizadas.

En consecuencia, mi postura profesional como docente en formación de la MAEM asume el compromiso de planear experiencias donde comprender, modelar, autorregular el propio aprendizaje y colaborar respetuosamente operen como dimensiones inseparables de una misma justicia educativa \citep{unesco2015replantear}.

\clearpage
\begin{thebibliography}{99}
\bibitem[Collin y Quigley(2021)]{eef2021feedback} Collin, J., y Quigley, A. (2021). \textit{Teacher feedback to improve pupil learning}. Education Endowment Foundation.
\bibitem[CONALEP(2025)]{conalep2025mecatronica} Colegio Nacional de Educación Profesional Técnica. (2025). \textit{Perfil de egreso: Profesional Técnico-Bachiller en Mecatrónica. Modelo Académico 2025}. Secretaría Académica, CONALEP.
\bibitem[Delors et~al.(1996)]{delors1996educacion} Delors, J., Al Mufti, I., Amagi, I., Carneiro, R., Chung, F., Geremek, B., Gorham, W., Kornhauser, A., Manley, M., Padrón Quero, M., Savané, M.-A., Singh, K., Stavenhagen, R., Suhr, M. W., y Nanzhao, Z. (1996). \textit{La educación encierra un tesoro: Informe a la UNESCO de la Comisión Internacional sobre la Educación para el Siglo XXI}. Santillana/UNESCO.
\bibitem[Pimienta Prieto(2012)]{pimienta2012estrategias} Pimienta Prieto, J. H. (2012). \textit{Estrategias de enseñanza-aprendizaje: Docencia universitaria basada en competencias}. Pearson Educación.
\bibitem[Sawyer(2008)]{sawyer2008optimizing} Sawyer, R. K. (2008). Optimizing learning: Implications of learning sciences research. En OECD, \textit{Innovating to learn, learning to innovate} (pp. 45--65). OECD Publishing.
\bibitem[Schunk(2012)]{schunk2012teorias} Schunk, D. H. (2012). \textit{Teorías del aprendizaje: Una perspectiva educativa} (6.ª ed.). Pearson Educación.
\bibitem[Scott(2015)]{scott2015futuro2} Scott, C. L. (2015). \textit{El futuro del aprendizaje 2: ¿Qué tipo de aprendizaje se necesita en el siglo XXI?} (Investigación y prospectiva en educación: contribuciones temáticas, núm. 14). UNESCO.
\bibitem[SEP(2023)]{sep2023pensamiento} Secretaría de Educación Pública. (2023). \textit{Programa de estudios del recurso sociocognitivo Pensamiento Matemático I}. Subsecretaría de Educación Media Superior, SEP.
\bibitem[UNESCO(2015)]{unesco2015replantear} UNESCO. (2015). \textit{Replantear la educación: ¿Hacia un bien común mundial?} Ediciones UNESCO.
\end{thebibliography}

\end{document}
